% Part: lambda-calculus
% Chapter: introduction
% Section: reduction

\documentclass[../../../include/open-logic-section]{subfiles}

\begin{document}

\olfileid{lam}{int}{red}
\olsection{اختزال حدود لامبدا}

ماذا يمكن أن نفعل بحدود لامبدا؟ نبسّطها. إذا كان~$M$ و$N$ أي حدين من
حدود لامبدا وكان~$x$ أي متغير، أمكننا استعمال~$\Subst{M}{N}{x}$ للدلالة
على نتيجة إحلال~$N$ محل~$x$ في~$M$، بعد إعادة تسمية أي متغيرات مربوطة
في~$M$ قد تتداخل، بعد الإحلال، مع المتغيرات الحرة في~$N$. فمثلًا،
\[
\Subst{(\lambd[w][xxw])}{yyz}{x} = \lambd[w][(yyz)(yyz)w].
\]

\begin{digress}
من الترميزات البديلة للإحلال~$[N/x]M$ و$[x/N]M$، وكذلك~$M[x/N]$.
فاحذر!
\end{digress}

من الوجهة الحدسية، لـ$(\lambd[x][M])N$ و$\Subst{M}{N}{x}$ المعنى نفسه؛
وتسمى عملية استبدال الحد الأول بالثاني \emph{تقلّص~$\beta$}.
ويسمى~$(\lambd[x][M])N$ \emph{موضع اختزال}، ويسمى
$\Subst{M}{N}{x}$ \emph{ناتج تقلّصه}. وبوجه عام، إذا أمكن تحويل حد~$P$
إلى~$P'$ بتقلّص~$\beta$ لحد جزئي ما، قلنا إن~$P$ \emph{يختزل بحسب
$\beta$ إلى~$P'$ في خطوة واحدة}، وكتبنا~$P \redone P'$. وإذا أمكننا
الحصول من~$P$ على~$P'$ بعدد ما من الاختزالات ذات الخطوة الواحدة (قد
يكون صفرًا)، فإن~$P$ \emph{يختزل بحسب~$\beta$} إلى~$P'$؛ ورمزيًا،
$P \red P'$. ويسمى الحد الذي لا يمكن اختزاله بحسب~$\beta$ أكثر
\emph{غير قابل للاختزال بحسب~$\beta$}، أو \emph{سويًا بحسب~$\beta$}.
وسنقول ``يختزل'' بدلًا من ``يختزل بحسب~$\beta$''، وما إلى ذلك، متى كان
السياق واضحًا.

لنتأمل بعض الأمثلة.
\begin{enumerate}
\item لدينا
\begin{align*}
(\lambd[x][xxy]) \lambd[z][z] & \redone (\lambd[z][z])(\lambd[z][z]) y \\
& \redone (\lambd[z][z]) y \\
& \redone y.
\end{align*}
\item قد يجعل ``تبسيط'' حد ما ذلك الحد أشد تعقيدًا:
\begin{align*}
(\lambd[x][xxy])(\lambd[x][xxy]) & \redone (\lambd[x][xxy])(\lambd[x][xxy])y \\
& \redone (\lambd[x][xxy])(\lambd[x][xxy])yy \\
& \redone \dots
\end{align*}
\item وقد يترك حدًا من دون تغيير أيضًا:
\[
(\lambd[x][xx])(\lambd[x][xx]) \redone (\lambd[x][xx])(\lambd[x][xx]).
\]
\item ويمكن كذلك اختزال بعض الحدود بأكثر من طريقة؛ فمثلًا،
\[
(\lambd[x][(\lambd[y][yx]) z)] v \redone (\lambd[y][yv]) z
\]
بتقليص التطبيق الأبعد إلى الخارج؛ و
\[
(\lambd[x][(\lambd[y][yx]) z)] v \redone (\lambd[x][zx]) v
\]
بتقليص التطبيق الأبعد إلى الداخل. لكن لاحظ أن كلا الحدين يختزل بعد ذلك،
في هذه الحالة، إلى الحد نفسه~$zv$.
\end{enumerate}

ليست النتيجة النهائية في المثال الأخير مصادفة، بل توضح خاصية عميقة
ومهمة لحساب لامبدا تعرف باسم ``خاصية تشيرش--روسر''.

\end{document}
